\documentclass[UTF8,11pt,a4paper,fontset=fandol]{ctexart}
\usepackage[margin=21mm,top=22mm,bottom=22mm]{geometry}
\usepackage{amsmath,amssymb,booktabs,array,tabularx,xcolor}
\usepackage{enumitem,fancyhdr,lastpage,hyperref,tcolorbox}
\hypersetup{colorlinks=true,urlcolor=blue!55!black,linkcolor=blue!55!black}
\definecolor{inkblue}{HTML}{355F86}
\definecolor{paleblue}{HTML}{EDF3F8}
\definecolor{palegreen}{HTML}{F0F5EF}
\setlength{\parindent}{0pt}
\setlength{\parskip}{5pt}
\linespread{1.18}
\setlist[itemize]{leftmargin=1.5em,itemsep=3pt,topsep=3pt}
\setlist[enumerate]{leftmargin=1.8em,itemsep=4pt,topsep=3pt}
\setlength{\abovedisplayskip}{7pt}
\setlength{\belowdisplayskip}{7pt}
\ctexset{section={format=\Large\bfseries\color{inkblue}},subsection={format=\normalsize\bfseries}}
\pagestyle{fancy}\fancyhf{}
\fancyhead[L]{\small CAHDRL：拉格朗日与比例--积分入门}
\fancyhead[R]{\small 按当前论文及实现整理}
\fancyfoot[C]{\small 第\thepage 页 / 共\pageref{LastPage}页}
\setlength{\headheight}{16pt}
\newtcolorbox{remember}{colback=paleblue,colframe=inkblue!45,boxrule=.5pt,arc=1mm,left=2mm,right=2mm,top=1mm,bottom=1mm}
\newtcolorbox{examplebox}{colback=palegreen,colframe=green!35!black!35,boxrule=.5pt,arc=1mm,left=2mm,right=2mm,top=1mm,bottom=1mm}
\newcommand{\lesson}[2]{\section*{#1\quad #2}}
\newcommand{\pos}[1]{\max\{0,#1\}}
\newcommand{\clamp}[2]{\min\{#1,\max\{0,#2\}\}}
\newcommand{\ind}{\mathbf{1}}
\begin{document}
\lesson{01}{先认识：你的算法为什么需要四个“调节旋钮”？}
\textbf{阅读目标：}看懂当前实现，不要求你先学会优化理论。本文不修改论文或代码；所有手算数字都明确是教学例子，不是实验结果。核对日期：2026年10月31日。

你希望调度算法\textbf{尽量省计算能耗}，但不能为了省电，把工作流拖得太晚，或让它的模型可靠性低于要求。因此，“省电”和“守约束”要一起考虑。

\begin{remember}
把乘子理解为“违反某类要求时，要扣多少分”的系数。\textbf{代价}表示这件事有多糟；\textbf{乘子}决定这类糟糕事情有多值得重视。两者相乘才是惩罚。
\end{remember}

你的两层策略做不同的事，所以有四个独立的系数：
\begin{center}
\renewcommand{\arraystretch}{1.45}
\begin{tabularx}{\linewidth}{@{}lXl@{}}
\toprule
乘子 & 对应的问题 & 作用位置\\\midrule
$\mu_D^H$ & 上层选服务器，要多重视整体Deadline代价？ & 上层动作评分\\
$\mu_R^H$ & 上层选服务器，要多重视可靠性代价？ & 上层动作评分\\
$\mu_D^L$ & 下层调频，要多重视降频增加的迟到？ & 下层训练奖励\\
$\mu_R^L$ & 下层调频，要多重视可靠性代价？ & 下层训练奖励\\\bottomrule
\end{tabularx}
\end{center}

\textbf{符号不要混：}
\begin{itemize}
\item 上标 $H$ 表示上层，$L$ 表示下层；下标 $D$ 表示Deadline，$R$ 表示可靠性。它们是标签，不是乘方。
\item $w$：一个工作流；$i$：一个任务；$e$：一整批工作流形成的训练轨迹。
\item $t$：批内的一次任务调度决策；$k$：控制器的一次有效更新。$k$不是任务编号。
\item $J$：控制器最近观测到的统计量；$b$：允许的预算；$J-b$：超预算多少。
\item 本文用 $U$ 表示“记住过去误差的累积量”。$U$与真正使用的乘子 $\mu$不是同一个变量。
\end{itemize}

\textbf{学习顺序：}先读第2--3页理解“为什么扣分”，再读第4--8页手算Deadline控制器，最后读可靠性与练习。不必先去背代码变量名。

\newpage
\lesson{02}{从约束问题，一步一步得到拉格朗日形式}
先看一个\textbf{教学用的标准问题}。$\theta$表示调度策略的参数：
\begin{align*}
&\text{目标：让平均能耗 } E(\theta) \text{ 尽量小；}\\
&\text{要求：代价 } C(\theta)\text{ 不超过预算 }b.
\end{align*}
把“不能超过”写成数学形式，再把右边移到左边：
\[
 C(\theta)\le b
 \quad\Longleftrightarrow\quad
 C(\theta)-b\le0.
\]
引入一个非负系数 $\lambda\ge0$，构造拉格朗日函数：
\begin{align*}
\mathcal L(\theta,\lambda)
 &= E(\theta)+\lambda\big(C(\theta)-b\big)\\
 &= E(\theta)+\lambda C(\theta)-\lambda b.
\end{align*}
\textbf{为什么是加号？}这里在\emph{最小化}。当 $C>b$ 时，多出来的正项会让这个方案变差。若改成最大化收益，令 $R(\theta)=-E(\theta)$，就变成：
\begin{align*}
-\mathcal L(\theta,\lambda)
 &=R(\theta)-\lambda C(\theta)+\lambda b.
\end{align*}
固定 $\lambda$ 时，$\lambda b$与动作、策略参数无关。因此，对策略来说，可以只比较 $R-\lambda C$。\textbf{但更新乘子时，预算 $b$不能丢。}

\subsection*{策略与乘子，更新方向为什么相反？}
对固定的乘子，策略尽量减小 $\mathcal L$；约束超标时，乘子往上调，给策略更大压力。标准对偶的定义是：
\[
g(\lambda)=\inf_\theta\mathcal L(\theta,\lambda),
\qquad \max_{\lambda\ge0}g(\lambda).
\]
不必急着求这个最优解。先注意：固定 $\theta$，$\mathcal L$对 $\lambda$的变化率就是 $C(\theta)-b$，所以常见的近似更新为：
\[
\lambda_{k+1}=\max\{0,\lambda_k+\alpha(C_k-b)\}.
\]
超标就增加，低于预算就减少。$\alpha$是每次调整的步幅。
\begin{remember}
这一页是标准理论的入门，不是说你的代码完整求解了这个对偶问题。你的实现使用不同层的代理代价、窗口统计、比例--积分和上限裁剪，应称为\textbf{拉格朗日启发的自适应惩罚}；不能由此宣称必然收敛或必然满足约束。理论出处见末页[1]。
\end{remember}

\newpage
\lesson{03}{四个乘子，究竟在哪儿扣分？}
\subsection*{上层：先算三个网络的值，再合成动作评分}
对候选服务器 $m$，三个网络分别预测：能耗收益 $Q_E(m)$、Deadline代价 $Q_D(m)$、可靠性代价 $Q_R(m)$。省略状态记号后：
\begin{align*}
S(m) &= Q_E(m)-\mu_D^H\max\{0,Q_D(m)\}-\mu_R^H Q_R(m),\\
m^* &=\arg\max_m S(m).
\end{align*}
$\arg\max$就是“选评分最高的服务器”。这里\textbf{Deadline预测做非负截断，可靠性预测没有同样截断}，按当前实现保留这一差别。

\begin{examplebox}
教学例子：可靠性预测暂设为0，只看Deadline系数。
\[
\begin{array}{c|cc}
&\text{A服务器}&\text{B服务器}\\\hline
Q_E&1.00&0.80\\
Q_D&0.20&0.02
\end{array}
\]
当 $\mu_D^H=0.5$：$S(A)=1-0.5\times0.2=0.90$，
$S(B)=0.8-0.5\times0.02=0.79$，选A。

当 $\mu_D^H=3$：$S(A)=1-3\times0.2=0.40$，
$S(B)=0.8-3\times0.02=0.74$，改选B。
\end{examplebox}
网络预测即使暂时不变，改变乘子也能改变选择。这里的预测值是训练尺度上的数，\textbf{不是直接拿J去减可靠性概率}。

\subsection*{下层：在SAC用于学习的奖励里扣分}
先把省电收益和局部时间代价处理成可比较的尺度：
\[
B_t=Z_E(\Delta E_i)-Z_h(h_t).
\]
$\Delta E_i$是相对满频节省的能量；$h_t$是降频消耗的局部时间裕量；$Z$表示代码中的标准化、尺度下限和裁剪。随后：
\[
r_t^L=B_t-\mu_D^L c_{D,t}^L-\mu_R^L c_{R,t}^L.
\]
若 $B_t=0.6,c_D^L=0.05,c_R^L=0.01,\mu_D^L=1,\mu_R^L=2$，则
$r_t^L=0.6-1\times0.05-2\times0.01=0.53$。

\textbf{不是直接“乘子变大就把频率改成某个值”。}乘子改变训练奖励，奖励影响Critic和Actor的学习，随后策略的调频行为才可能改变。神经网络学习并不保证每一次动作都按同一方向变化。

\newpage
\lesson{04}{上层Deadline：先算清楚“整批失败”}
你的系统有两个不同的90\%，需要分两层读：
\begin{enumerate}
\item $\Gamma=0.90$：\textbf{一批内部}至少90\%的工作流按时完成。
\item $\eta=0.90$：希望\textbf{这样的批次}达标的概率至少90\%。
\end{enumerate}
对第 $e$ 批，先统计按时完成的工作流：
\[
\mathrm{SR}_e=\frac{\text{该批按时完成数}}{\text{该批工作流总数}}.
\]
再把它变成一个“这批是否失败”的0/1标记：
\[
z_e=\begin{cases}
1,&\mathrm{SR}_e<0.90,\\
0,&\mathrm{SR}_e\ge0.90.
\end{cases}
\]
例如训练一批20个工作流，18个按时完成：$18/20=0.90$，$z_e=0$；17个按时完成：$17/20=0.85$，$z_e=1$。测试一批50个时，45个按时完成也算这批达标。

设窗口里已有 $n$ 批，最多保留最近60批：
\begin{align*}
J_H &=\frac{z_{e-n+1}+z_{e-n+2}+\cdots+z_e}{n},\\
b_H&=1-\eta=1-0.90=0.10,\\
\text{当前误差 }\varepsilon_H&=J_H-b_H=J_H-0.10.
\end{align*}
\textbf{这里 $J_H$是失败批次比例，不是平均工作流迟到比例。}

\begin{examplebox}
最近60批中12批失败：$J_H=12/60=0.20$。
\[
\varepsilon_H=0.20-0.10=0.10>0.
\]
需要提高Deadline惩罚。若只有3批失败：
$J_H=3/60=0.05$，误差 $=-0.05$，当前压力可以减小。
\end{examplebox}

上层前50个训练episode不更新Deadline乘子；之后也要有足够观测，当前实现至少需要10批。\textbf{60是窗口容量，不是必须等满60批才能开始。}上述比例只是有限样本统计，不是总体成功率的保证。

\newpage
\lesson{05}{比例--积分：把PI完全写开}
PI是\textbf{比例--积分}（Proportional--Integral）：
\begin{itemize}
\item 比例项P：只看\textbf{现在}超标多少，立刻响应。
\item 积分项I：把\textbf{过去每次}的超标量逐步累积，保留记忆。
\end{itemize}
下面每个 $k$都表示一次\textbf{允许执行的更新}。设当前统计量为 $J_k$、预算为 $b$、积分步幅为 $\tau$、比例系数为 $K_P$、乘子上限为 $M$。你的代码依次做：
\begin{align*}
\text{第1步：误差 }\varepsilon_k &= J_k-b,\\
\text{第2步：积分候选 }\widetilde U_{k+1} &= U_k+\tau(J_k-b),\\
\text{第3步：积分裁剪 }U_{k+1} &=\min\{M,\max\{0,\widetilde U_{k+1}\}\},\\
\text{第4步：比例项 }P_k &=K_P(J_k-b),\\
\text{第5步：输出候选 }\widetilde\mu_{k+1} &=U_{k+1}+K_P(J_k-b),\\
\text{第6步：乘子裁剪 }\mu_{k+1} &=\min\{M,\max\{0,\widetilde\mu_{k+1}\}\}.
\end{align*}
先把数限制到不小于0，再限制到不大于 $M$。例如 $\min\{6,\max\{0,-0.3\}\}=0$；输入7则输出6。\textbf{比例项本身可以为负；不是每一步都截成正数。}

为了看见“积分记忆”，\textbf{暂时假设两处裁剪都没有触发}，且 $U_0=1$：
\begin{align*}
U_1&=1+\tau\varepsilon_0,\\
U_2&=1+\tau\varepsilon_0+\tau\varepsilon_1,\\
U_3&=1+\tau\varepsilon_0+\tau\varepsilon_1+\tau\varepsilon_2,\\
\mu_{k+1}&=1+\tau(\varepsilon_0+\varepsilon_1+\cdots+\varepsilon_k)
               +K_P\varepsilon_k.
\end{align*}
\begin{remember}
上面的累加展开\textbf{仅在没有触发裁剪时成立}；实际计算始终按前面的六步。积分状态 $U$存过去，比例项只取当前误差。\textbf{下次累加用旧 $U$，不是旧乘子 $\mu$。}PI没有微分项，不应叫完整PID。相关控制视角见末页[2]。
\end{remember}

\newpage
\lesson{06}{上层Deadline：把你的参数全部代进去}
上层参数是：预算0.10、积分步幅0.10、比例系数3、上限6；初始 $U_H=\mu_D^H=1$。每次有效更新完整写成：
\begin{align*}
U_{H,k+1}&=\min\{6,\max\{0,U_{H,k}+0.10(J_{H,k}-0.10)\}\},\\
\mu_{D,k+1}^H&=\min\{6,\max\{0,U_{H,k+1}+3(J_{H,k}-0.10)\}\}.
\end{align*}
\textbf{第一次：}设已经过预热，窗口60批中12批失败，更新前积分状态为1：
\begin{align*}
J_{H,0}&=12/60=0.20,\\
\varepsilon_{H,0}&=0.20-0.10=0.10,\\
U_{H,1}&=\min\{6,\max\{0,1+0.10\times0.10\}\}=1.01,\\
P_0&=3\times0.10=0.30,\\
\mu_{D,1}^H&=\min\{6,\max\{0,1.01+0.30\}\}=1.31.
\end{align*}
\textbf{第二次：}下一次窗口中11批失败，例如最旧的失败批次移出，新来的批次达标。\textbf{接着用刚才的 $U_{H,1}=1.01$：}
\begin{align*}
J_{H,1}&=11/60\approx0.183333,\\
\varepsilon_{H,1}&=11/60-0.10\approx0.083333,\\
U_{H,2}&=1.01+0.10\times0.083333\approx1.018333,\\
P_1&=3\times0.083333=0.25,\\
\mu_{D,2}^H&=1.018333+0.25=1.268333.
\end{align*}
\begin{remember}
两次都超标，\textbf{积分状态确实增加了}；但当前误差变小，比例项减少，所以实际乘子可以从1.31降至1.268333。\textbf{“只要超标，乘子就一定比上次大”对PI并不成立。}
\end{remember}
常见错误是第二次算成 $1.31+0.10\times0.083333+0.25$。这把旧比例项也当成积分状态，重复累加了。正确起点是1.01。

\newpage
\lesson{07}{下层Deadline：先算降频究竟多造成多少迟到}
下层不是简单地把“自己迟到”记为1，而是看当前固定调度中，\textbf{降频额外增加的归一化迟到}。本页不要求你先学完整的责任分配算法。

对工作流 $v$，设相对Deadline为 $D_v$，实际响应时间为 $M_v$。先定义归一化迟到量：
\[
\ell_v(M_v)=\max\{0,M_v/D_v-1\}.
\]
固定任务放置、服务器执行顺序和原始分配时刻，把整批任务都改成满频，重算响应时间 $M_v^{\mathrm{all}}$。下层观测代价是：
\[
d_v^L=\min\{1,\max\{0,\ell_v(M_v)-\ell_v(M_v^{\mathrm{all}})\}\}.
\]
\begin{examplebox}
Deadline为100秒，实际130秒，固定调度改满频后110秒：
\begin{align*}
\ell_v(M_v)&=\max\{0,130/100-1\}=0.30,\\
\ell_v(M_v^{\mathrm{all}})&=\max\{0,110/100-1\}=0.10,\\
d_v^L&=\min\{1,\max\{0,0.30-0.10\}\}=0.20.
\end{align*}
满频也会迟到的那部分，没有在这里全部算到调频头上。
\end{examplebox}
一批有 $N_e$个工作流，先求这批的平均影响：
\[
a_e=\frac{d_1^L+d_2^L+\cdots+d_{N_e}^L}{N_e}.
\]
最近最多40个 $a_e$按从大到小排列，取最大的25\%再平均。若已有 $n$批，取 $q=\max\{1,\lceil0.25n\rceil\}$ 个；$\lceil\ \rceil$表示向上取整：
\[
J_L=\frac{a_{(1)}+\cdots+a_{(q)}}q,\qquad a_{(1)}\ge\cdots\ge a_{(n)}.
\]
满40批时就是最坏10批的均值。教学用4批 $0.01,0.03,0.06,0.12$，最坏25\%只取0.12，所以 $J_L=0.12$，不是四批平均值0.055。

\textbf{实际实现是有限窗口的“最坏25\%均值”估计。}它是尾部风险代理，不等同于总体分布上的精确CVaR，也不是ESP。

\newpage
\lesson{08}{下层Deadline：把PI及归零的情况写开}
下层参数：预算0.08、积分步幅0.20、比例系数25、上限8。初始 $U_L=\mu_D^L=1$。每次有效更新：
\begin{align*}
U_{L,k+1}&=\min\{8,\max\{0,U_{L,k}+0.20(J_{L,k}-0.08)\}\},\\
\mu_{D,k+1}^L&=\min\{8,\max\{0,U_{L,k+1}+25(J_{L,k}-0.08)\}\}.
\end{align*}
\textbf{例A：}已经过预热，$U_{L,0}=1,J_{L,0}=0.12$。
\begin{align*}
\varepsilon_{L,0}&=0.12-0.08=0.04,\\
U_{L,1}&=1+0.20\times0.04=1.008,\\
P_0&=25\times0.04=1,\\
\mu_{D,1}^L&=\min\{8,\max\{0,1.008+1\}\}=2.008.
\end{align*}
\textbf{例B：独立的另一种情形}，更新前也设 $U_L=1$，但 $J_L=0$：
\begin{align*}
\varepsilon_L&=0-0.08=-0.08,\\
U_L^{\mathrm{new}}&=1+0.20\times(-0.08)=0.984,\\
P&=25\times(-0.08)=-2,\\
\widetilde\mu_D^L&=0.984-2=-1.016,\\
\mu_D^{L,\mathrm{new}}&=\min\{8,\max\{0,-1.016\}\}=0.
\end{align*}
\begin{remember}
预算0.08\textbf{不表示允许8\%的工作流迟到}，也不表示Deadline缩短8\%。它限制的是上一页定义的尾部统计量。下层系数变成0可能是这个更新规则的直接结果，不一定是程序故障。
\end{remember}
当 $\mu_D^L=0$，下层奖励中的这一项 $\mu_D^L c_D^L$为0，但局部时间项 $Z_h(h_t)$、上层Deadline代价、可靠性项并不因此全部消失。也不能仅凭归零就断言Deadline模块无用或有用。

下层前20个训练episode不更新Deadline乘子；之后至少要有4个有效批次统计。40仍只是窗口容量。重新加载模型时控制器可能恢复已保存状态，不能总把实际 $U$当成1；这里的1只是从头训练的初值。

\newpage
\lesson{09}{上层可靠性：非负缺口，越界才增加}
用 $R_w$表示工作流可靠性要求，用 $\mathcal R_w$表示实际模型可靠性。当前要求通常为 $R_w=0.90$。上层看的不是一个简单的0/1，而是\textbf{对数可靠性缺口}：
\[
d_w^R=\max\{0,\ln R_w-\ln\mathcal R_w\}.
\]
$\ln$是自然对数，可以用计算器的ln键。由于ln单调递增，实际值低于要求时，$\ln R_w-\ln\mathcal R_w$为正；满足要求时截为0。

为什么可靠性会出现对数？若独立任务的可靠性相乘得到工作流可靠性：
\[
\mathcal R_w=r_1r_2\cdots r_n,
\qquad
\ln\mathcal R_w=\ln r_1+\ln r_2+\cdots+\ln r_n.
\]
乘积转成求和，便于累计和计算；这不代表独立性假设已经被现实数据验证。

上层可靠性没有这里介绍的额外比例项。按每个完成工作流的缺口更新一次，步幅0.001、上限100：
\[
\mu_R^{H,\mathrm{new}}=
\min\{100,\max\{0,\mu_R^{H,\mathrm{old}}
+0.001\max\{0,\ln R_w-\ln\mathcal R_w\}\}\}.
\]
训练收集完成后可以批量遍历这些工作流执行更新；\textbf{不是每个任务都更新一次}。

\begin{examplebox}
若旧系数为1、要求0.90、实际0.85：
\begin{align*}
d_w^R&=\ln(0.90/0.85)\approx0.0571584,\\
\mu_R^{H,\mathrm{new}}&=1+0.001\times0.0571584\approx1.00005716.
\end{align*}
若下一个工作流实际可靠性为0.95：缺口
$\max\{0,\ln(0.90/0.95)\}=0$，系数\textbf{保持不变}，不是往下减。
\end{examplebox}

\textbf{结论：}在这条规则下，上层可靠性乘子只增不减，达到上限后不能继续增大。“一切都达标就四个乘子一起下降”是不对的。

\newpage
\lesson{10}{下层可靠性：看最近工作流的0/1失败标记}
下层可靠性控制器改看每个已完成工作流是否达到要求：
\[
v_w=\begin{cases}
1,&\mathcal R_w<R_w,\\
0,&\mathcal R_w\ge R_w.
\end{cases}
\]
对最近最多100个工作流求均值。窗口中已有 $n_R$个时：
\[
J_R=\frac{v_1+v_2+\cdots+v_{n_R}}{n_R}.
\]
预算 $b_R=0.005$、步幅0.5、上限50；这条也没有额外P项：
\[
\mu_R^{L,\mathrm{new}}=
\min\{50,\max\{0,\mu_R^{L,\mathrm{old}}+0.5(J_R-0.005)\}\}.
\]
\textbf{满100个工作流的例子：}
\begin{align*}
\text{1个失败： }J_R&=1/100=0.01,\\
\mu_R^{L,\mathrm{new}}&=1+0.5(0.01-0.005)=1.0025;\\[2pt]
\text{全都达标： }J_R&=0/100=0,\\
\mu_R^{L,\mathrm{new}}&=1+0.5(0-0.005)=0.9975.
\end{align*}
上面是两种独立情形，旧系数都取1。$0.005=0.5\%$，不是5\%；窗口满100时，一个失败就是1\%，已经高于0.5\%预算。\textbf{训练允许的小预算不等于测试表中把Rel.=100\%改成99.5\%。}

\subsection*{务必分开：调乘子的统计量，与奖励里的代价}
\begin{itemize}
\item 此处 $J_R$用最终0/1标记，告诉控制器最近是否经常违反要求。
\item SAC奖励里的 $c_{R,t}^L$仍是降频额外引起的\emph{投影可靠性缺口}，不是把每个任务都简单罚1。
\end{itemize}
若投影可靠性在频率 $f$下为 $\widehat{\mathcal R}_{w,t}(f)$，定义
\begin{align*}
u_t(f)&=\frac{\max\{0,\ln R_w-\ln\widehat{\mathcal R}_{w,t}(f)\}}{-\ln R_w},\\
c_{R,t}^L&=\max\{0,u_t(f_i)-u_t(F_{m_i})\}.
\end{align*}
先算当前频率的缺口，再减去同一任务满频时的缺口；负差值截成0。这就是为什么“奖励代价”和“控制器观测”不能混写成一个量。

\newpage
\lesson{11}{把整套训练流程串起来，再自己算一遍}
\begin{enumerate}
\item 用当前 $\mu_D^H,\mu_R^H$给候选服务器打分，选择任务放置。
\item SAC给出频率，环境实际执行，记录能耗、时刻和可靠性信息。
\item 一批结束，形成上层整体代价和下层额外代价；完成对应的反馈分配。
\item 上层可靠性按工作流缺口更新；神经网络用各自数据和损失学习。下层学习奖励使用\emph{当前}乘子。
\item 轨迹结束时，满足门槛后更新上下层Deadline控制器及下层可靠性控制器；新的乘子供后续训练使用。实现中的各更新不是同时求出一个解析最优解。
\item 测试时冻结网络权重和这些乘子，不利用测试结果继续调它们。
\end{enumerate}

\subsection*{自测题：先挡住答案，用计算器算}
\begin{enumerate}[label=\arabic*.]
\item 上层Deadline：已经允许更新，$U_H=0.5,J_H=0.15$。新积分量、比例项、乘子各是多少？
\item 下层Deadline：已经允许更新，$U_L=0.8,J_L=0.04$。为什么乘子可能变成0？
\item 上层可靠性：$\mu_R^H=2$、实际可靠性等于要求0.90。更新后是多少？
\item 下层可靠性：$\mu_R^L=1$，最近100个工作流有2个失败。更新后是多少？
\end{enumerate}

\subsection*{逐步答案}
\textbf{1.} 误差 $0.15-0.10=0.05$；
积分 $0.5+0.10\times0.05=0.505$；
比例项 $3\times0.05=0.15$；输出 $0.505+0.15=\boxed{0.655}$。

\textbf{2.} 误差 $0.04-0.08=-0.04$；
积分 $0.8+0.20\times(-0.04)=0.792$；
比例项 $25\times(-0.04)=-1$；候选值 $0.792-1=-0.208$；裁剪后 $\boxed{0}$。

\textbf{3.} 缺口 $\max\{0,\ln0.90-\ln0.90\}=0$；输出 $2+0.001\times0=\boxed{2}$。

\textbf{4.} 失败比例 $2/100=0.02$；误差 $0.02-0.005=0.015$；输出 $1+0.5\times0.015=\boxed{1.0075}$。

\begin{remember}
能独立做对这四道题，就已经理解了四个乘子各自看什么、怎么调、能不能下降。之后再看源码，不必从神经网络内部开始找。
\end{remember}

\newpage
\lesson{12}{最后一页：参数速查、代码位置和结论边界}
\begingroup\small
\linespread{1.02}\selectfont
\setlength{\parskip}{3pt}
\setlist[itemize]{itemsep=2pt,topsep=2pt}
\renewcommand{\arraystretch}{1.15}
\begin{center}
\begin{tabularx}{\linewidth}{@{}lXXXX@{}}
\toprule
 & 上层Deadline & 下层Deadline & 上层可靠性 & 下层可靠性\\\midrule
观测 & 失败批次比例 & 额外迟到尾均值 & 对数缺口 & 失败工作流比例\\
预算 & 0.10 & 0.08 & 0 & 0.005\\
积分/调整步幅 & 0.10 & 0.20 & 0.001 & 0.5\\
额外比例系数 & 3 & 25 & 无 & 无\\
乘子上限 & 6 & 8 & 100 & 50\\
初始乘子 & 1 & 1 & 1 & 1\\
统计窗口 & 最近60批 & 最近40批最坏25\% & 逐工作流 & 最近100工作流\\
预热 & 前50轮 & 前20轮 & 无此门槛 & 无额外预热\\\bottomrule
\end{tabularx}
\end{center}
Deadline的两个积分状态也从1开始。可靠性下层窗口非空即可统计，不必等满100个；以上均针对当前论文完整模型，不混用Raw或预算0.01的诊断版本。

\subsection*{代码导航：看函数，不必通读整仓库}
\begin{itemize}
\item 上层Deadline：\path{src/D3QN/deadline_lagrangian.py}。\\
函数：{\small\path{DeadlineLagrangianController.update}}。
\item 下层Deadline：\path{src/sac/deadline_lagrangian.py}。\\
函数：{\small\path{LowDeadlineLagrangianController.update}}。
\item 上层可靠性：\path{src/D3QN/reliability_aware_high.py}。\\
函数：{\small\path{ReliabilityLagrangian.update}}。
\item 下层可靠性：\path{src/sac/reliability_lagrangian.py}。\\
函数：{\small\path{LowReliabilityLagrangianController.update}}。
\item 当前耦合训练入口：\\
{\footnotesize\path{tools/train_coupled_hierarchical_d3qn_sac.py}}。\\
函数{\small\path{install_low_dual_attribution}}将下层可靠性控制器的输入替换为最终工作流0/1标记；不要只读底层收集器就误认输入。
\item 调用时机与下层Deadline统计构造：\\
{\footnotesize\path{tools/train_hierarchical_d3qn.py}}。\\
看轨迹末更新及函数{\footnotesize\path{_workflow_counterfactual_low_deadline_credit}}。
\end{itemize}

\textbf{三句最重要的话：}乘子不是硬约束开关；不同预算对应不同统计量；有限样本达标不等于总体可行性证明。你的论文当前用“拉格朗日启发”是对实现边界的准确说明。

\subsection*{可继续阅读的原始资料}
\footnotesize
[1] Boyd、Vandenberghe，\emph{Convex Optimization}，第5章Duality。标准拉格朗日、对偶函数与非负乘子的定义：
\url{https://see.stanford.edu/materials/lsocoee364a/05Duality.pdf}

[2] Stooke、Achiam、Abbeel，\emph{Responsive Safety in Reinforcement Learning by PID Lagrangian Methods}，ICML 2020。用于理解乘子更新的控制视角，不意味着你的PI实现完整复现了其PID算法：
\url{https://proceedings.mlr.press/v119/stooke20a.html}

\textbf{本项目依据：}当前\path{paper/manuscript/main2.tex}的约束反馈与乘子更新部分，以及上述控制器和\path{configs/v3_sm_74_legacy_train.json}。讲义中的小数演算仅为教学，未改变模型、训练或实验数据。
\endgroup
\end{document}
